\documentclass[11pt]{article}
\usepackage[margin=1in]{geometry}
\usepackage{amsmath,amssymb}
\newcommand{\finalanswer}[1]{\par\medskip\noindent\fbox{\parbox{0.94\linewidth}{\textbf{Answer:} #1}}}
\title{STAT 412 Homework 4 --- Question 13}
\author{}
\date{}
\begin{document}
\maketitle

\section*{Problem}
A 100-watt bulb has mean life $\mu=900$ hours with standard deviation
$\sigma=60$ hours. At most, what percentage of bulbs fail to last even 660
hours? Assume the distribution is symmetric about the mean.

\section*{Setup}
Chebyshev's theorem bounds the \emph{two-sided} tail:
$P(|X-\mu|\ge k\sigma)\le \dfrac{1}{k^2}$. ``Failing to last 660 hours'' means
$X\le 660$, which is the \emph{lower} tail. Because the distribution is symmetric
about the mean, each tail carries half of the two-sided bound:
$P(X\le \mu-k\sigma)\le \dfrac{1}{2k^2}$.

\section*{Work}
First express $660$ in standard deviations below the mean:
\[
k=\frac{\mu-660}{\sigma}=\frac{900-660}{60}=\frac{240}{60}=4 .
\]
Then
\[
P(X\le 660)\le \frac{1}{2k^2}=\frac{1}{2(4)^2}=\frac{1}{32}=0.03125=3.125\% .
\]

\finalanswer{At most $\dfrac{1}{32}=3.125\%$ of the bulbs fail to last 660 hours.}
\end{document}
